\chapter{Teorie -- Dráhy družic}
\label{4-teorie-drahy-druzic}

\section{Pohyb družice (Keplerovský pohyb)}

Pohyb družice lze idealizovat jako pohyb v~gravitačním poli Země, kde je Země uvažována jako koule se středově souměrným rozložením hmoty a nejsou uvažovány žádné vnější vlivy. Takto zjednodušený model pohybu se označuje jako Keplerovský pohyb a představuje základní teoretický popis oběžných drah~\cite{bursa_kostelecky_1994}. Gravitační pole ideální Země lze popsat gravitačním potenciálem

\begin{equation}
V~= -\frac{GM}{\rho},
\end{equation}
kde:
\begin{itemize}
    \item $G$ je gravitační konstanta,
    \item $M$ hmotnost Země,
    \item $\rho = |\mathbf{r}|$ vzdálenost družice od středu Země,
    \item $\mathbf r$ je polohový vektor družice vzhledem ke středu Země.
\end{itemize}
Zrychlení družice v~potenciálovém poli je dáno záporným gradientem potenciálu

\begin{equation}
\ddot{\mathbf r} = - \nabla V.
\end{equation}
Pro potenciál gravitačního pole ideální Země pak platí

\begin{equation}
\nabla V~= \frac{GM}{\rho^3}\mathbf r,
\end{equation}
a výsledná rovnice pohybu má tvar

\begin{equation}
\ddot{\mathbf r} = -\frac{GM}{\rho^3}\mathbf r.
\end{equation}

Tuto rovnici lze rozložit jako soustavu tří diferenciálních rovnic druhého řádu, která při zadaných počátečních podmínkách (poloze a rychlosti), jednoznačně určuje pohyb družice. Z~jejího tvaru je zřejmé, že zrychlení je v~každém okamžiku orientováno do středu Země, a pohyb tedy probíhá v~centrálním silovém poli. V~centrálním silovém poli je zachován moment hybnosti družice

\begin{equation}
\mathbf h = \mathbf r \times \dot{\mathbf r},
\end{equation}
z~čehož plyne, že pohyb probíhá v~jedné rovině a že plošná rychlost průvodiče je konstantní. Tato vlastnost odpovídá druhému Keplerovu zákonu~\cite{bursa_kostelecky_1994}.

Řešením rovnice pohybu v~centrálním gravitačním poli jsou kuželosečky. V~případě umělých družic Země se jedná především o~eliptické dráhy, které lze popsat vztahem

\begin{equation}
\rho = \frac{a(1 - e^2)}{1 + e \cos \nu},
\end{equation}

kde 
\begin{itemize}
    \item $a$ je hlavní poloosa, 
    \item $e$ je excentricita,
    \item $\nu$ je pravá anomálie.
\end{itemize}

Keplerovy zákony popisují geometrické vlastnosti dráhy i časový průběh pohybu. V~tomto smyslu nepředstavují samostatný předpoklad, ale důsledek pohybu v~centrálním gravitačním poli.

Polohu a orientaci eliptické dráhy lze jednoznačně popsat pomocí šesti Keplerovských elementů: hlavní poloosy $a$, excentricity $e$, sklonu dráhy $i$, rektascenze vzestupného uzlu $\Omega$, argumentu perigea $\omega$ a střední anomálie $M$. Tyto parametry představují alternativní popis dráhy vůči kartézské poloze a rychlosti a v~případě Keplerovského pohybu jsou konstantní v~čase.

Keplerovský pohyb představuje ideální řešení problému dvou těles. Ve skutečnosti je však pohyb družice ovlivněn řadou dalších sil, které vedou k~odchylkám od ideální eliptické dráhy a ke změnám Keplerovských elementů v~čase. Tím vzniká potřeba uvažovat obecnější, rušený pohyb družice.

\section{Rušený pohyb družice}

Ve skutečnosti se družice nepohybují v~ideálním centrálním gravitačním poli, protože Země není homogenní koule a na družice současně působí další gravitační i negravi\-tační vlivy. Tyto vlivy způsobují odchylky od ideálního Keplerovského pohybu a takový pohyb se označuje jako rušený pohyb družice~\cite{bursa_kostelecky_1994}. Pohyb družice proto nelze popsat pouze Keplerovským modelem, ale je nutné uvažovat obecnější rovnici pohybu ve tvaru

\begin{equation}
\ddot{\mathbf{r}} = -\frac{GM}{\rho^3}\mathbf{r} + \mathbf{a}_{\text{pert}},
\end{equation}
kde 
\begin{itemize}
    \item $\mathbf{a}_{\text{pert}}$ představuje součet poruchových zrychlení.
\end{itemize}
Poruchová zrychlení \(\mathbf{a_{pert}}\) lze rozdělit na gravitační a negravitační. 

\paragraph{Mezi gravitační poruchy patří:}
\begin{itemize}
\item nesférické gravitační pole Země popsané harmonickým rozvojem,
\item časově proměnné slapové deformace Země,
\item gravitační působení Slunce, Měsíce a planet.
\end{itemize}

\paragraph{Mezi negravitační poruchy patří:}
\begin{itemize}
\item tlak slunečního záření,
\item záření Země (albedo a infračervené vyzařování),
\item atmosférický odpor (významný zejména pro nízké oběžné dráhy).
\end{itemize}

\section{Přesné určení drah družic}

Přesné určení drah družic (POD) je proces, při němž se na základě měření a dynamického modelu pohybu odhaduje dráha družic. Dynamický model zahrnuje nejen ideální centrální gravitační pole Země, ale i jeho odchylky a další gravitační i negravitační rušivé vlivy.  Odhad dráhy je prováděn nepřímo prostřednictvím odhadu parametrů tohoto modelu, přičemž samotná trajektorie je výsledkem řešení rovnice pohybu. Cílem je určit polohu a rychlost družice v~čase s~co nejvyšší přesností. Výsledkem tohoto procesu jsou produkty přesných drah, typicky poskytované ve formátu SP3, které obsahují diskrétní hodnoty polohy a případně rychlosti družice v~čase.

V~případě systému DORIS je určování dráhy založeno na měření Dopplerova posunu signálu vysílaného pozemními stanicemi a přijímaného na palubě družice. Měřený Dopplerův posun souvisí se změnou vzdálenosti mezi družicí a stanicí v~čase, tedy s~relativní rychlostí ve směru spojnice družice--stanice. Z~časového průběhu těchto měření během průletů družice nad stanicemi lze následně odhadovat parametry  dynamického modelu tak, aby výsledná dráha co nejlépe odpovídala pozorováním~\cite{ids_what_is_doris}.

\section{Dynamická parametrizace dráhy}

Dynamická parametrizace dráhy je způsob modelování drah družic, při kterém je dráha popsána jako řešení rovnice pohybu s~uvažováním gravitačních i negravitačních sil. Na rozdíl od čistě geometrických metod je tento přístup založen na fyzikálním modelu pohybu družice. Dráha družice je obvykle popisována na časově omezených úsecích, tzv. obloucích, v~jejichž rámci jsou aplikovány fyzikální modely a současně odhadovány parametry popisující její pohyb. Tyto parametry určují konkrétní řešení rovnice pohybu. Princip spočívá v~tom, že odhadované parametry jsou voleny tak, aby výsledná trajektorie co nejlépe odpovídala vstupním datům, například diskrétním souřadnicím družice ve formátu SP3~\cite{bernese_5_2}. V~rámci jednoho oblouku jsou aplikovány zejména následující fyzikální modely:

\begin{itemize}

\item \textbf{Gravitační modely}
\begin{itemize}
\item Nesférické gravitační pole Země, popsané pomocí harmonického rozvoje
\item Slapové jevy
\item Gravitační působení Slunce, Měsíce a blízkých planet
\end{itemize}

\item \textbf{Modely negravitačních sil}
\begin{itemize}
\item Apriorní model hustoty atmosféry MSIS-86
\item Apriorní model tlaku slunečního záření
\item Záření odražené a vyzařované Zemí (albedo a infračervené záření)
\end{itemize}

\end{itemize}

\noindent Současně jsou v~rámci jednotlivých oblouků odhadovány následující parametry:

\begin{itemize}

\item \textbf{Počáteční podmínky}
\begin{itemize}
\item Šest parametrů určujících stav družice v~epoše na začátku oblouku. Tyto parametry lze vyjádřit jako polohový vektor a vektor rychlosti, nebo ekvivalentně šesti Keplerovskými elementy.
\end{itemize}

\item \textbf{Odhadované parametry negravitačních sil}
\begin{itemize}
\item Měřítkový koeficient atmosférického odporu
\item Měřítkový koeficient tlaku slunečního záření
\end{itemize}

\item \textbf{Empirické parametry}
Jsou parametry, které nemají jasný fyzikální význam, ale zachycují nemodelované nebo nepřesně modelované vlivy. Mají periodu přibližně odpovídající oběžné době družice. Tyto parametry jsou obvykle definovány ve směrech:
\begin{itemize}
\item R -- radiální (ve směru polohového vektoru),
\item T -- transverzální (ve směru pohybu družice),
\item N -- normálový (kolmo na rovinu dráhy),
\end{itemize}
a jsou vyjádřeny pomocí amplitudy a fáze (sinusové a kosinusové členy). Jejich úlohou je absorbovat zbytkové chyby modelování~\cite{stepanek_2023_gop_itrf2020}.

\end{itemize}

\section{Interpolace drah}

Kromě dynamické parametrizace dráhy z~diskrétních bodů lze pro určení polohy družice v~konkrétní epoše využít také interpolaci mezi známými body, která je čistě matematickým přístupem a nevyžaduje zahrnutí fyzikálních vlivů do výpočtu. Pohyb družice po oběžné dráze je nelineární, a proto je při požadavku na vysokou přesnost vhodné použít interpolační metodu, která kromě polohy zohledňuje také první derivaci polohy, tedy rychlost družice~\cite{zeitlhoefler_2024_interpolation}.

\paragraph{Hermitova interpolace}\leavevmode\\
Hermitova interpolace v~uzlech respektuje nejen hodnotu funkce $f(t)$, ale i její první derivaci $f'(t)$. V~případě drah družic představuje funkce $f(t)$ polohu družice $\mathbf r(t)$ a její derivace $f'(t)$ odpovídá rychlosti družice $\mathbf v(t)=\dot{\mathbf r}(t)$. Jsou dány uzly \(t_0, t_1, \dots, t_n\), v~nichž jsou známy hodnoty $f(t_k)$ a $f'(t_k)$. Hermitův interpolační polynom má stupeň nejvýše $2n+1$~\cite{zeitlhoefler_2024_interpolation}. V~obecném tvaru lze Hermitovu interpolaci zapsat jako:

\begin{equation}
H_{2n+1}(t) = \sum_{k=0}^{n} f(t_k)\,H_{n,k}(t) + \sum_{k=0}^{n} f'(t_k)\,\hat H_{n,k}(t).
\label{eq:hermite_general}
\end{equation}

\paragraph{Lagrangeovy bázové funkce}\leavevmode\\
Základem konstrukce Hermitova interpolačního polynomu jsou Lagrangeovy bázové funkce:

\begin{equation}
L_{n,k}(t)=
\prod_{\substack{i=0 \\ i\neq k}}^{n}
\frac{t-t_i}{t_k-t_i}.
\label{eq:lagrange_basis}
\end{equation}
Pro sestavení Hermitových bázových funkcí je dále potřeba derivace Lagrangeovy bázové funkce v~jejím vlastním uzlu, která je dána vztahem

\begin{equation}
L'_{n,k}(t_k)=
\sum_{\substack{m=0 \\ m\neq k}}^{n}
\frac{1}{t_k-t_m}.
\label{eq:lagrange_deriv}
\end{equation}

\paragraph{Hermitovy báze}\leavevmode\\
Pro danou epochu $t^\ast$

\begin{equation}
\Delta t_k = t^\ast - t_k.
\end{equation}
Hermitovy bázové funkce mají tvar

\begin{equation}
H_{n,k}(t^\ast)
=
\left(
1 - 2 \Delta t_k L'_{n,k}(t_k)
\right)
L_{n,k}^2(t^\ast),
\end{equation}

\begin{equation}
\hat H_{n,k}(t^\ast)
=
\Delta t_k
L_{n,k}^2(t^\ast).
\end{equation}
Funkce $H_{n,k}$ váží příspěvek polohy v~uzlu $t_k$,
zatímco $\hat H_{n,k}$ váží příspěvek rychlosti.

\paragraph{Výsledný tvar}

Interpolovaná poloha družice v~epoše $t^\ast$ je určena vztahem

\begin{equation}
\mathbf r(t^\ast)
=
\sum_{k=0}^{n}
\mathbf r(t_k)\, H_{n,k}(t^\ast)
+
\sum_{k=0}^{n}
\mathbf v(t_k)\, \hat H_{n,k}(t^\ast).
\end{equation}
kde:
\begin{itemize}
    \item \(\mathbf{r}\) je polohový vektor \((x, y, z)\)
    \item \(\mathbf{v}\) je vektor rychlosti \((v_x, v_y, v_z)\)
\end{itemize}

Interpolace je prováděna složkově pro jednotlivé kartézské souřadnice, přičemž bázové funkce jsou společné pro všechny složky vektoru. Interpolační polynom je obvykle konstruován na základě uzlů z~okolí požadované epochy.

\section{Transformace rozdílů do systému RTN}

Pro interpretaci rozdílů drah se často používá lokální orbitální souřadnicový systém RTN~\cite{zeitlhoefler_2024_interpolation}. Ten tvoří tři na sebe kolmé složky:

\begin{itemize}
\item R – radiální směr (směr polohového vektoru, od středu Země k~družici),
\item T – transverzální směr (ve směru pohybu, kolmo na R v~rovině dráhy),
\item N – normálový směr (kolmý k~rovině dráhy).
\end{itemize}

\noindent Jednotkové vektory lze definovat jako:
\begin{equation}
\hat{\mathbf{e}}_R = \frac{\mathbf{r}}{|\mathbf{r}|},
\end{equation}

\begin{equation}
\hat{\mathbf{e}}_N = \frac{\mathbf{r} \times \mathbf{v}}{|\mathbf{r} \times \mathbf{v}|},
\end{equation}

\begin{equation}
\hat{\mathbf{e}}_T = \hat{\mathbf{e}}_N \times \hat{\mathbf{e}}_R.
\end{equation}

\noindent Rozdíl dvou drah $\Delta \mathbf{r}$ lze rozložit jako:

\begin{equation}
\Delta r_R = \Delta \mathbf{r} \cdot \hat{\mathbf{e}}_R,
\end{equation}

\begin{equation}
\Delta r_T = \Delta \mathbf{r} \cdot \hat{\mathbf{e}}_T,
\end{equation}

\begin{equation}
\Delta r_N = \Delta \mathbf{r} \cdot \hat{\mathbf{e}}_N.
\end{equation}
Radiální složka je obvykle nejcitlivější na chyby modelování gravitačního pole,
zatímco transverzální složka je citlivá na chyby v~časové synchronizaci.

\section{Metriky hodnocení rozdílů drah}

Pro kvantifikaci rozdílů mezi dvěma řešeními se používá:

\begin{itemize}
\item okamžitý rozdíl v~RTN,
\item 3D rozdíl:
\begin{equation}
|\Delta \mathbf{r}| = \sqrt{\Delta X^2 + \Delta Y^2 + \Delta Z^2},
\end{equation}
\item střední kvadratická chyba (RMS):
\begin{equation}
\mathrm{RMS} = \sqrt{\frac{1}{n}\sum_{i=1}^{n} (\Delta r_i)^2},
\end{equation}
\item střední kvadratická chyba po odečtení průměru (RMS0):
\begin{equation}
\mathrm{RMS0} = \sqrt{\frac{1}{n}\sum_{i=1}^{n} (\Delta r_i - \overline{\Delta r})^2},
\end{equation}
kde $\overline{\Delta r} = \frac{1}{n}\sum_{i=1}^{n} \Delta r_i$ je průměrná hodnota rozdílů.
\end{itemize}
Zatímco RMS zahrnuje i případný systematický posun, RMS0 jej odečítá
a vyjadřuje pouze rozptyl rozdílů kolem průměru. Tyto metriky umožňují
objektivní porovnání řešení různých analytických center a posouzení
vlivu použité metodiky.